\documentclass[10pt,a4paper]{amsart}
\usepackage[margin=19mm]{geometry}
\usepackage{amsmath,amssymb,mathtools}
\usepackage[scr=rsfs,cal=euler]{mathalfa}
\usepackage{xfrac}
\usepackage[unicode,hidelinks,bookmarks=false]{hyperref}
\newcommand{\ie}{i.e.\ }
\newcommand{\cf}{cf.\ }
\newcommand{\resp}{respectively\ }
\newcommand{\Subsection}{\subsection}
\providecommand{\nobreakdash}{\nobreak\hbox{-}}
\setlength{\emergencystretch}{2em}
\multlinegap0pt
\usepackage{bm}
\usepackage{bbm}
\usepackage{dsfont}
\usepackage{xspace}
\usepackage{cancel}
\usepackage{mathtools}
\usepackage{amsthm}
\makeatletter
\@ifundefined{lemma}{\newtheorem{lemma}{Lemma}}{}
\@ifundefined{theorem}{\newtheorem{theorem}{Theorem}}{}
\makeatother
\newcommand{\I}{\mathcal{I}_N}
\newcommand{\K}{\mathcal{K}}
\newcommand{\R}{\mathbb{R}}
\newcommand{\dd}{\mathrm{d}}
\newcommand{\eps}{\varepsilon}
\newcommand{\pa}{\partial}
\title[Convergence of a semi-explicit scheme for a one dimensional ]{Convergence of a semi-explicit scheme for a one dimensional periodic nonlocal eikonal equation modeling dislocation dynamics}
\author{Diana Al Zareef, Ahmad El Hajj, Hassan Ibrahim, Antoine Zurek}
\date{Source: arXiv:2501.18428v2 (CC BY 4.0)}
\begin{document}
\maketitle

\noindent\emph{Source and licence.} Convergence of a semi-explicit scheme for a one dimensional periodic nonlocal eikonal equation modeling dislocation dynamics. Version arXiv:2501.18428v2 (CC BY 4.0), distributed under the Creative Commons Attribution 4.0 International licence, as recorded in the arXiv licence metadata for this exact version and shown on the abstract page https://arxiv.org/abs/2501.18428v2. Full rights notice: licenses/arxiv-2501.18428-CC-BY-4.0.txt.\par\smallskip

\noindent\textit{Editorial scope.} Counted unit: the Lemma whose source label is \texttt{lem.fpartial} in the article (source0.tex lines 1037--1042, source label \texttt{lem.fpartial}), delivered together with the article's own complete original proof of it (source0.tex lines 1044--1109, one \texttt{proof} environment of 66 lines). No mathematics is written, repaired or re-derived: statement and proof are the article's own text line for line. Required context retained uncounted: def.LP (lines 185--187); 1.initperiod (lines 374--376); lem.init.perio.to.nonperiod (lines 380--386); estimationLinP (lines 390--392); estimationL1P (lines 394--396); schemeIC (lines 442--444); def.dis.grad (lines 453--455); h(x) (lines 465--467); 1.deff (lines 489--494); thm.scheme (lines 496--522); Linftybounddis (lines 510--512); TVestim (lines 514--516); gradentdis (lines 518--520); defzeta (lines 871--873); lem.boundinitent (lines 930--939); lem.LinftyL1estime (lines 969--974); def.tau (lines 996--998); decompo.patuPeps (lines 1000--1002); tau in theta (lines 1012--1014); estimationLPPLP (lines 1027--1028). Every definition, display and label of those lines is retained unchanged. Omitted from this package and not redistributed: 1--184, 188--373, 377--379, 387--389, 393--393, 397--441, 445--452, 456--464, 468--488, 495--495, 513--513, 517--517, 521--870, 874--929, 940--968, 975--995, 999--999, 1003--1011, 1015--1026, 1029--1036, 1043--1043, 1110--1534. Nothing inside a retained excerpt is omitted or altered: every retained body is delivered exactly as the article's lines are written, so no omission receipt of any kind accompanies this package. Citations: the retained excerpts carry no citation command at all, so no bibliography entry of the article is transcribed and no reference is redistributed. Reference adaptation: none was needed and none was made; every reference inside the delivered unit resolves inside it, so no unresolved reference or citation is produced. Printed slips kept literal: no printed word, symbol, hypothesis or number of the counted statement or of its proof was altered. Layout and class: the article's own class is \texttt{amsart} with packages including a4wide, enumerate, enumitem, amsmath, color, inputenc, geometry, amssymb, amsthm, graphicx, subcaption, cite, bigints, breqn; the wrapper is the lane's pinned \texttt{amsart} editorial wrapper at 10pt on A4, which restates the theorem-like environments the retained text needs and adds the article's own macro definitions that the retained text needs (\texttt{\textbackslash I}, \texttt{\textbackslash K}, \texttt{\textbackslash R}, \texttt{\textbackslash dd}, \texttt{\textbackslash eps}, \texttt{\textbackslash pa}), with extra wrapper packages (bm, bbm, dsfont, xspace, cancel, mathtools). The delivered numbering is the unit's own. Assets: no retained span contains a figure environment, a graphics-inclusion command, a PDF-page inclusion, a table or a TikZ picture; no image file is copied into this package, the package declares no asset dependency and no image file is redistributed with it. Licence: version arXiv:2501.18428v2 of this article is distributed under the Creative Commons Attribution 4.0 International licence, as recorded in the arXiv licence metadata for this exact version and shown on the abstract page https://arxiv.org/abs/2501.18428v2. A full rights notice is delivered at licenses/arxiv-2501.18428-CC-BY-4.0.txt. Characters: every retained excerpt is pure ASCII, so the delivered engine, XeLaTeX with its default Latin Modern text font, reports no missing character.
\begin{equation}\label{def.LP}
L^P=\frac{v_0(P)-v_0(-P)}{2P}, 
\end{equation} 

\begin{align}\label{1.initperiod}
    u^P_0(x) = v_0(x) - L^P \, x, \quad \mbox{for } x \in [-P,P],
\end{align}

\begin{lemma}\label{lem.init.perio.to.nonperiod}
Let assumption~\textbf{\emph{(H1)}} holds and assume that $P\ge 1$. Then, the function $u^P_0$, given by~\eqref{1.initperiod}, satisfies, uniformly with respect to $P$,
\begin{align*}
    u^P_0 \in L^\infty(I_P), \quad \pa_x u^P_0 \in L^1(I_P)\cap L\,\log \,L(I_P).
\end{align*}    
Moreover, the continuous function $x \in [-P,P)  \mapsto u^P_0(x) + L^Px$ is nondecreasing.
\end{lemma}

\begin{align}\label{estimationLinP}
    \|u^P_0\|_{L^\infty(I_P)}\leq \|v_0\|_{L^\infty(\R)} + \frac12 \left|v_0(P)-v_0(-P)\right|\leq 2\|v_0\|_{L^\infty(\R)},
\end{align}

\begin{align}\label{estimationL1P}
    \|\pa_x u^P_0\|_{L^1(I_P)} \leq \|\pa_x v_0\|_{L^1(\R)} + 2\|v_0\|_{L^\infty(\R)} \le  4\|v_0\|_{L^\infty(\R)}. 
\end{align}

\begin{align}\label{schemeIC}
u_i^{P, 0} =  u^P_0(x_i), \quad \forall i \in \I.
\end{align}

\begin{align}\label{def.dis.grad}
    \theta^{P,n}_{i+1/2} := \dfrac{u^{P,n}_{i+1}-u^{P,n}_i}{\Delta x}, \quad \forall i\in\I.
\end{align}

  \begin{equation}\label{h(x)}
  \lambda^M_i \left[u^{P, n + 1}\right]= \sum_{j \in \I} \Delta x \, \sigma_{M,j}^P \, u^{P, n + 1}_{i - j}, \quad \forall i \in \I, 
 \end{equation}

\begin{equation}\label{1.deff}
    0 \leq f(x) = \begin{cases}
        x \ln(x) + \frac{1}{e} & \quad\mbox{if } x\geq 1/e,\\
        0  & \quad \mbox{if } 0 \leq x \leq  1/e.
    \end{cases}          
\end{equation}

\begin{theorem}[Properties of the scheme]\label{thm.scheme}
Let assumptions~\textbf{\emph{(H2)}}--\textbf{\emph{(H4)}} hold, and assume that
\begin{equation}\label{delta t}
    \Delta t < \frac{1}{10 (L^P) \|\mathcal{K}\|_{L^1 (\mathbb{R})}}\,\min\left\{\frac{1}{\|u^P_0\|_{L^\infty(I_P)}\, e^{10 (L^P) T ||\mathcal{K}||_{L^1(\mathbb{R})}} + 1}, 1\right\},
\end{equation}
and
\begin{equation}\label{delta t over delta x}
   \frac{\Delta t}{\Delta x} < \frac{1}{10 \|\mathcal{K}\|_{L^1 (\mathbb{R})}}
    \min\left\{ 
\frac{1}{\left(\|u^P_0\|_{L^\infty(I_P)}\, e^{10 (L^P)T \|\mathcal{K}\|_{L^1(\mathbb{R})}} + 1\right)}, 
\frac{1}{3 \|u^P_0\|_{L^\infty(I_P)}\, e^{10 (L^P)T \|\mathcal{K}\|_{L^1(\mathbb{R})}}} 
    \right\}.
\end{equation}
Then, for any $1\leq n\leq N_T$, the scheme~\eqref{schemeIC}--\eqref{h(x)} admits a unique solution $u^{P,n} = (u^{P,n}_i)_{i\in\I}$ such that
\begin{align}\label{Linftybounddis}
 \max_{ i \in \I} \, |u_{i}^{P, n}| & \leq \|u^P_0\|_{L^\infty(I_P)} \, e^{10 (L^P) T \|\mathcal{K}\|_{L^1(\mathbb{R})}}, 
\end{align} 
and $\theta_{i + 1/2}^{P, n} + L^P$ is nonnegative for any $i\in\I$ where we recall definition~\eqref{def.dis.grad} of $\theta^{P,n}_{i+1/2}$. Furthermore, this solution satisfies the following discrete total variation estimate:
\begin{align}\label{TVestim}
 TV(u^{P, n} )= \sum_{i \in \I} \left| u^{P,n}_{i+1}-u^{P,n}_i\right| \leq TV(u^P_0 )\, e^{10 (L^P) T \|\mathcal{K}\|_{L^1(\mathbb{R})}}, \quad \forall 1\leq n \leq N_T.
 \end{align}
Finally, there exists a positive constant $\zeta_0$ only depending on $T$ and $\|\mathcal{K}\|_{L^1(\mathbb{R})}$ (see~\eqref{defzeta}), such that the following discrete gradient entropy estimate holds
 \begin{align}\label{gradentdis}
 \sum_{i \in\I} \Delta x \, f\left(\theta_{i + 1/2}^{P, n} + L^P\right) \leq \sum_{i \in \I} \Delta x \, f\left(\theta_{i + 1/2}^{P, 0}+ L^P\right) + \zeta_0 \, TV(u^P_0), \quad \mbox{for } 1\leq n \leq N_T,
\end{align}
where we recall definition~\eqref{1.deff} of the function $f$.
\end{theorem}

\begin{align}\label{defzeta}
    \zeta \leq 5T \, \|\K\|_{L^1(\R)} \, e^{10 T \,\|v_0\|_{L^\infty(\R)} \, \|\K\|_{L^1(\R)}} \, \left(\frac{1}{e \ln(2)}+\|v_0\|_{L^\infty(\R)}\right) :=  \zeta_0.
\end{align}

\begin{lemma}\label{lem.boundinitent}
Let the assumptions of Theorem~\ref{thm.scheme} hold. Then, there exists a constant $C_0>0$ 
only depending on $v_0$ 
 such that
\begin{align*}
    I_0=\sum_{i\in\I} \Delta x \, f\left(\theta^{P,0}_{i+1/2}+L^P\right) \leq C_0, 
\end{align*}
where $L^P$ is defined in \eqref{def.LP}. 

\end{lemma}

\begin{lemma}\label{lem.LinftyL1estime}
Let the assumptions of Theorem~\ref{thm.scheme} hold. Then, there exists a constant $C_1>0$ only depending on $v_0$, $T$ and $\K$ such that
\begin{align}\label{LinftyL1.paxu.patu}
    \|\pa_x u^{P,\eps}_M\|_{L^\infty(0,T;L^1(I_P))} + \|\pa_t u^{P,\eps}_M\|_{L^\infty(0,T;L^1(I_P))} \leq C_1.
\end{align}
\end{lemma}

\begin{align}\label{def.tau}
\tau_i^{P, n +1/2} := \frac{u_i^{P, n + 1} - u_i^{P, n}}{\Delta t}, \quad \forall i \in \I, \, 0\leq n \leq N_T-1. 
\end{align}

\begin{equation}\label{decompo.patuPeps}
\partial_t u^{P, \eps}_M(x,t) = \bigg(\frac{x - x_i}{\Delta x}\bigg) \, \tau_{i + 1}^{P, n +1/2}+\bigg( 1 - \frac{x - x_i}{\Delta x}\bigg)\, \tau_i^{P, n +1/2}. 
\end{equation}

\begin{equation}\label{tau in theta}
\tau_i^{P, n +1/2} = \lambda_i[u^{P, n + 1}]_+ \, \left(\theta_{i +1/2}^{P,n} + L^P\right) - \lambda_i[u^{P,n + 1}]_- \, \left(\theta_{i-1/2}^{P, n} + L^P\right).
\end{equation}

\begin{equation}\label{estimationLPPLP}L^P  \leq \|v_0\|_{L^\infty(\R)},  
\quad \mbox{and}\quad PL^P  \leq \|v_0\|_{L^\infty(\R)} \end{equation}

\begin{lemma}\label{lem.fpartial}
Let the assumptions of Theorem~\ref{thm.scheme} hold. Then, there exists a constant $C_2>0$ only depending on  $v_0$, $T$ and $\K$ such that
\begin{align}
    \int_{I_P} f\left(\pa_x u^{P,\eps}_M(x,t)+L^P \right) \, \dd x + \int_{I_P} f\left(\left|\pa_t u^{P,\eps}_M(x,t)\right| \right) \, \dd x \leq C_2, \quad \mbox{for a.e. }t\in(0,T).
\end{align}
\end{lemma}

\begin{proof}
First, for $(x,t)\in(x_i,x_{i+1})\times[t_n,t_{n+1}]$ with $i\in \I$ and $0\le n \le N_T-1$, we write
\begin{equation}\label{uepsilonx + L}
\partial_x u^{P, \eps}_M(x,t)  + L^P=  \bigg(\frac{t - t_n}{\Delta t}\bigg) \, \left(\theta_{i +1/2}^{P, n + 1} + L^P\right) + \bigg(1 - \frac{t - t_n}{\Delta t}\bigg) \, \left(\theta_{i +1/2}^{P, n }+ L^P\right). 
\end{equation}
Now, using the convexity of $f$, we obtain
\[
f\left(\partial_x u^{P, \eps}_M(x,t)+L^P \right) \leq \bigg(\frac{t - t_n}{\Delta t}\bigg) \, f\left(\theta_{i +1/2} ^{P, n + 1} + L^P\right)+ \bigg(1 - \frac{t - t_n}{\Delta t}\bigg) \, f\left(\theta_{i +1/2}^{P, n }+ L^P\right).
\]
Therefore, by integrating in space, we have
\begin{multline*}
\int_{I_P} f\left(\partial_x u^{P, \eps}_M(x,t)+L^P \right)\, \dd x \leq \bigg(\frac{t - t_n}{\Delta t}\bigg) \, \sum_{i\in\I} \Delta x\, f\left(\theta^{P,n+1}_{i+1/2}+L^P\right)\\
+ \bigg(1 - \frac{t - t_n}{\Delta t}\bigg) \, \sum_{i\in\I} \Delta x \, f\left(\theta^{P,n}_{i+1/2}+L^P\right).
\end{multline*}
Thanks to estimate~\eqref{gradentdis} we get
\begin{align*}
\int_{I_P} f\left(\partial_x u^{P, \eps}_M(x,t)+L^P \right)\, \dd x &\leq \sum_{i\in\I} \Delta x \, f\left(\theta^{P,0}_{i+1/2}+L^P\right) + \zeta_0 TV(u^P_0), 
\end{align*}
where $\zeta_0$ is defined~\eqref{defzeta}. Applying Lemma \ref{lem.boundinitent}, we deduce that 
\begin{align*}
   \int_{I_P} f\left(\partial_x u^{P, \eps}_M(x,t)+L^P \right)\, \dd x 
   \leq C_0 + \zeta_0  TV(u^P_0)
   \le C_0 + 4\zeta_0  \|v_0\|_{L^\infty(\R)}
\end{align*}
where for the last inequality we have used Lemma~ \ref{lem.init.perio.to.nonperiod} (cf. \eqref{estimationL1P}).\\
 Let us now establish the estimate for $\pa_t u^{P,\eps}_M$. For this purpose, using~\eqref{decompo.patuPeps}, we notice that for all $(x,t) \in [x_i,x_{i+1}]\times(t_n,t_{n+1})$
\begin{align*}
      \left|\partial_t u^{P, \eps}_M(x,t)\right| &\leq \bigg(\frac{x - x_i}{\Delta x}\bigg) \,\left|\tau_{i + 1}^{P, n + 1/2}\right| + \bigg( 1 - \frac{x - x_i}{\Delta x}\bigg) \, \left|\tau_{i } ^{P, n +1/2}\right|,
\end{align*}
where we recall definition~\eqref{def.tau} of $\tau^{P,n+1/2}_i$. Hence, the convexity of $f$ yields
\begin{multline*}
    f\left(\left|\partial_t u^{p, \eps}_M(x,t)\right|\right) \leq \bigg(\frac{x - x_i}{\Delta x}\bigg)\,f\left(\left|\tau_{i + 1}^{P, n + 1/2}\right|\right) + \bigg( 1 - \frac{x - x_i}{\Delta x}\bigg) \, f\left(\left|\tau_{i } ^{P, n +1/2}\right|\right)\\
    \leq f\left(\left|\tau_{i + 1}^{P, n + 1/2}\right|\right) +  f\left(\left|\tau_{i } ^{P, n +1/2}\right|\right).
\end{multline*}
Let us for instance consider the second term in the right hand side. Then, according to the relation~\eqref{tau in theta}, we have
\begin{align*}
    \left|\tau_i^{P, n +1/2}\right| \leq \lambda_i[u^{P, n + 1}]_+ \, \left(\theta_{i +1/2}^{P,n} + L^P\right) + \lambda_i[u^{P,n + 1}]_- \, \left(\theta_{i-1/2}^{P, n} + L^P\right).
\end{align*}
Now, let us assume that $\lambda_i[u^{P,n+1}]\ge 0$. This yields
\begin{align*}
    \left|\tau_i^{P, n +1/2}\right| \leq |\lambda_i[u^{P, n + 1}]| \, \left(\theta_{i +1/2}^{P,n} + L^P\right).
\end{align*}
Then, we notice thanks to the $L^\infty$ estimate~\eqref{Linftybounddis} and the discrete total variation estimate~\eqref{TVestim}, that it holds
\begin{align*}
    \left|\tau_i^{P, n +1/2}\right| \leq \Lambda\, \left(\theta_{i +1/2}^{P,n} + L^P\right),
\end{align*}
where
\begin{align*}
    \Lambda := \max\left(1,5\|u^P_0\|_{L^\infty(I_P)} \, \|\K\|_{L^1(\R)} \, e^{10 (L^P)T \, \|\K\|_{L^1(\R)}}\right).
\end{align*}
Hence, using the elementary estimate $f(\gamma \theta) \leq \gamma f(\theta)+\theta f(\gamma)$ for all $\gamma\geq 1$ and $\theta \geq 0$, we have
\begin{align*}
    f\left(\left|\tau_i^{P, n +1/2}\right|\right) \leq \Lambda\, f\left(\theta_{i +1/2}^{P,n} + L^P\right) + \left(\theta_{i +1/2}^{P,n} + L^P\right)  \, f\left(\Lambda\right).
\end{align*}
Of course, one obtain a similar bound if $\lambda_i[u^{P,n+1}]\le 0$, so that
\begin{align}\label{bound.ftau}
    f\left(\left|\tau_i^{P, n +1/2}\right|\right) &\leq \sum_{\pm} \left[\Lambda\, f\left(\theta_{i \pm 1/2}^{P,n} + L^P\right) + \left(\theta_{i \pm 1/2}^{P,n} + L^P\right)  \, f\left(\Lambda\right) \right].
\end{align}
Thereby, we end up with
\begin{align*}
  I_1=  \int_{I_P} &f\left(\left|\pa_t u^{P,\eps}_M(x,t)\right|\right) \, \dd x \leq \sum_{i\in\I} \Delta x \, f\left(\left|\tau_{i + 1}^{P, n + 1/2}\right|\right) + \sum_{i\in\I} \Delta x\,  f\left(\left|\tau_{i } ^{P, n +1/2}\right|\right)\\
    &\leq 4\Lambda \left(C_0 + \zeta_0 \,TV(u^P_0)\right) + 8 PL^P \, f(\Lambda) + 4f(\Lambda) \, TV(u^P_0) \, e^{10 (L^P)T \, \|K\|_{L^1(\R)}},
\end{align*}
where we have used inequality~\eqref{bound.ftau}, the discrete gradient entropy inequality~\eqref{gradentdis} togther with Lemma~\ref{lem.boundinitent}, and the discrete total variation estimate~\eqref{TVestim}.\\
To complete the proof, we proceed as in the proof of Lemma \ref{lem.LinftyL1estime}. Indeed, we use \eqref{estimationLPPLP}, Lemma \ref{lem.init.perio.to.nonperiod} \eqref{estimationLinP}-\eqref{estimationL1P} and Lemma~\ref{lem.boundinitent},  in order to bound from above $I_1$ by a constant  $C_2$  independent of $M$ and $P$. 
\end{proof}

\end{document}
